%% ============================================================
%%  Chapter 2 — Background and Related Work
%% ============================================================
\chapter{Background and Related Work}
\label{ch:background}

This chapter surveys the theoretical and empirical foundations of HoneyIQ.
Section~\ref{sec:bg_rl} introduces reinforcement learning and Markov Decision
Processes.
Section~\ref{sec:bg_dqn} covers Deep Q-Networks.
Section~\ref{sec:bg_markov} presents Discrete-Time Markov Chains.
Section~\ref{sec:bg_killchain} describes the Lockheed Martin Cyber Kill Chain.
Section~\ref{sec:bg_honeypot} reviews honeypot technology.
Section~\ref{sec:bg_unsw} introduces the UNSW-NB15 dataset.
Section~\ref{sec:bg_rf} covers Random Forests.
Section~\ref{sec:bg_interpretable} discusses interpretable machine learning.
Section~\ref{sec:bg_related} reviews related work.
Section~\ref{sec:bg_repctrl} introduces reputation-based access control and
bounded control, the two design references behind the dynamic response
mechanisms developed in Chapter~\ref{ch:methodology}.

%% ── 2.1 Reinforcement Learning ───────────────────────────────
\section{Reinforcement Learning}
\label{sec:bg_rl}

Reinforcement learning (RL) is a computational framework in which an
\emph{agent} learns to make decisions through interaction with an
\emph{environment}~\citep{sutton2018reinforcement}.
Unlike supervised learning, the agent receives no labelled examples; instead,
it receives a scalar \emph{reward} signal $r_t$ after each action, and its
objective is to discover a \emph{policy} $\pi$ that maximises the expected
cumulative discounted return:
%
\begin{equation}
  G_t = \sum_{k=0}^{\infty} \gamma^k \cdot r_{t+k+1},
  \label{eq:return}
\end{equation}
%
where $\gamma \in [0, 1)$ is the \emph{discount factor} controlling the
relative weight of immediate versus future rewards.
In HoneyIQ, $\gamma = 0.99$, reflecting the importance of long-term campaign
containment over myopic step-by-step responses.

\subsection{Markov Decision Process Formalisation}
\label{sec:mdp}

The interaction is formalised as a Markov Decision Process (MDP)
$\langle \mathcal{S}, \mathcal{A}, P, R, \gamma \rangle$:
%
\begin{itemize}
  \item $\mathcal{S}$ --- the \emph{state space}.
        In HoneyIQ this is a continuous 24-dimensional observation space
        combining one-hot attack type, one-hot kill chain stage, threat level,
        attack count, escalation rate, and attacker intent.
  \item $\mathcal{A}$ --- the \emph{action space}: five discrete honeypot
        responses (ALLOW, LOG, TROLL, BLOCK, ALERT).
  \item $P(s' \mid s, a)$ --- the \emph{transition dynamics}.
        In HoneyIQ, the transition is determined by the attacker's Markov
        chain, which advances independently of the defender's action.
        This reflects the assumption that the attacker cannot observe the
        defender's specific response in real time.
  \item $R(s, a, s')$ --- the \emph{reward function} encoding domain
        knowledge about correct honeypot behaviour (Section~\ref{sec:reward}).
  \item $\gamma = 0.99$ --- discount factor.
\end{itemize}

The Markov property states that the future is conditionally independent of
the past given the present state:
$P(s_{t+1} \mid s_t, a_t, \ldots, s_0, a_0) = P(s_{t+1} \mid s_t, a_t)$.
This holds exactly in HoneyIQ because the 24-dimensional state vector
encodes all information relevant to predicting the next state.

\subsection{Value Functions and the Bellman Equation}

The \emph{state-value function} $V^\pi(s)$ gives the expected return from
state $s$ under policy $\pi$:
%
\begin{equation}
  V^\pi(s) = \mathbb{E}_\pi\!\left[\sum_{k=0}^{\infty} \gamma^k r_{t+k+1}
                                   \,\Big|\, s_t = s\right].
\end{equation}
%
The \emph{action-value function} $Q^\pi(s, a)$ extends this to
state-action pairs:
%
\begin{equation}
  Q^\pi(s, a) = \mathbb{E}_\pi\!\left[\sum_{k=0}^{\infty} \gamma^k r_{t+k+1}
                                      \,\Big|\, s_t = s,\, a_t = a\right].
\end{equation}
%
The optimal action-value function $Q^*$ satisfies the Bellman optimality
equation:
%
\begin{equation}
  Q^*(s, a) = \mathbb{E}\!\left[r + \gamma \cdot
              \max_{a'} Q^*(s', a') \,\Big|\, s, a\right].
  \label{eq:bellman}
\end{equation}
%
The optimal policy is $\pi^*(s) = \arg\max_a Q^*(s, a)$.

\subsection{Q-Learning}

Q-learning~\citep{watkins1992q} estimates $Q^*$ through incremental updates
from observed transitions $(s, a, r, s')$:
%
\begin{equation}
  Q(s, a) \leftarrow Q(s, a) + \alpha\!\left[
      r + \gamma \cdot \max_{a'} Q(s', a') - Q(s, a)
  \right],
\end{equation}
%
where $\alpha \in (0, 1]$ is the learning rate.
Under mild technical conditions, Q-learning converges to $Q^*$ in the tabular
setting, but the table grows exponentially with the size of the state space,
making it infeasible for continuous or high-dimensional observations.

%% ── 2.2 Deep Q-Networks ──────────────────────────────────────
\section{Deep Q-Networks (DQN)}
\label{sec:bg_dqn}

DQN~\citep{mnih2015humanlevel} resolves the scalability limitation by
approximating $Q(s, a; \theta)$ with a neural network parameterised by
$\theta$.
The seminal work demonstrated human-level performance on 49 Atari games using
raw pixel inputs, establishing DQN as a practical algorithm for large state
spaces.

\subsection{Neural Network Approximator}

In HoneyIQ, the policy network is a three-hidden-layer multi-layer perceptron:
%
\begin{equation}
  \text{Input}(24)
  \;\xrightarrow{\text{Linear+LayerNorm+ReLU}}\; 256
  \;\xrightarrow{\text{Linear+LayerNorm+ReLU}}\; 128
  \;\xrightarrow{\text{Linear+LayerNorm+ReLU}}\; 64
  \;\xrightarrow{\text{Linear}}\; \mathbb{R}^5.
\end{equation}
%
The network receives a 24-dimensional state vector and outputs five Q-values,
one per honeypot action.
Layer normalisation~\citep{ba2016layer} is applied after each linear layer;
it normalises activations across features rather than across the batch,
making it appropriate for small replay-buffer mini-batches and environments
where the input distribution shifts during training.

\subsection{Experience Replay}

Consecutive environment transitions are highly correlated, which violates
the independent and identically distributed (i.i.d.) assumption that
underlies stochastic gradient descent.
DQN addresses this by storing transitions $(s, a, r, s', \text{done})$ in
a \emph{replay buffer} (circular deque of capacity 15\,000) and sampling
uniformly at random for each gradient update.
This decorrelates the training distribution and allows each transition to
contribute to multiple updates.

\subsection{Target Network}

If the same network provides both the Q-value predictions and the
Bellman targets, the targets shift with every gradient step, leading
to oscillation and divergence.
DQN introduces a \emph{target network} $Q(s, a; \theta^-)$ whose parameters
are periodically synchronised from the policy network:
%
\begin{equation}
  y_i = r_i + \gamma \cdot \max_{a'} Q(s'_i, a'; \theta^-).
  \label{eq:td_target}
\end{equation}
%
In HoneyIQ, $\theta^-$ is hard-copied from the policy network every
150 steps, providing target stability over a window that spans multiple
episodes given the 200-step episode length.

\subsection{Huber Loss and Gradient Clipping}

The temporal-difference (TD) error is minimised with the Huber (SmoothL1)
loss:
%
\begin{equation}
  \mathcal{L}(\delta) =
  \begin{cases}
    \tfrac{1}{2}\delta^2 & |\delta| \le 1 \\
    |\delta| - \tfrac{1}{2} & \text{otherwise,}
  \end{cases}
  \quad \delta = y_i - Q(s_i, a_i; \theta).
\end{equation}
%
Huber loss is quadratic near zero (matching MSE for small errors) and linear
for large errors (making it less sensitive to outlier Q-values than MSE).
Gradients are additionally clipped to a maximum $\ell_2$ norm of 10.0 to
prevent explosive updates in early training.

\subsection{Epsilon-Greedy Exploration}

DQN uses an $\varepsilon$-greedy policy: with probability $\varepsilon$
the agent takes a random action (exploration); otherwise it selects the
action with the highest estimated Q-value (exploitation).
$\varepsilon$ is annealed exponentially:
%
\begin{equation}
  \varepsilon_{t+1} \leftarrow \max\!\left(\varepsilon_{\min},\;
                                            \varepsilon_t \cdot d\right),
  \quad d = 0.997,\quad \varepsilon_{\min} = 0.05.
\end{equation}
%
Starting from $\varepsilon = 1.0$, the agent explores almost uniformly for
the first $\approx$\,100 episodes before transitioning to predominantly
exploitative behaviour.
This warm-up period is essential: the replay buffer must accumulate diverse
transitions before the Q-value surface is well-defined.

\subsection{Limitations of DQN in Security Contexts}

Despite its empirical success, DQN presents two significant challenges for
operational cybersecurity deployment.

First, \textbf{opacity}: the Q-values produced by a multi-layer perceptron
are not directly interpretable by a human analyst.
When a DQN agent chooses BLOCK over LOG for a given traffic pattern, it is
impossible to trace the decision to a human-readable rule without
post-hoc explanation techniques~\citep{lundberg2017unified}, which introduce
their own approximation errors.

Second, \textbf{training instability}: the interaction of function
approximation, bootstrapping (the TD target depends on the same function
being updated), and off-policy data can produce oscillating or diverging
Q-value estimates~\citep{van2016deep}.
While the stabilising mechanisms (replay buffer, target network, gradient
clipping) substantially mitigate this in practice, convergence is not
guaranteed and must be verified empirically.

These limitations motivate the SEDM as an alternative policy design
(Section~\ref{sec:sedm}), and are examined empirically --- not merely cited
--- in Section~5.2 and Section~5.3, where the DQN's own recorded training
behaviour is used as direct evidence when deciding how the dynamic-adjustment
mechanisms introduced in \S\ref{sec:dynamic_response} should be designed.

%% ── 2.3 Markov Chains ────────────────────────────────────────
\section{Markov Chains}
\label{sec:bg_markov}

A Discrete-Time Markov Chain (DTMC)~\citep{norris1998markov} is a stochastic
process $\{X_t\}_{t \geq 0}$ over a finite state space $\mathcal{X}$ that
satisfies the Markov property:
%
\begin{equation}
  P(X_{t+1} = j \mid X_t = i,\, X_{t-1} = i_{t-1},\, \ldots)
  = P(X_{t+1} = j \mid X_t = i) \eqqcolon T_{ij}.
\end{equation}
%
The $|\mathcal{X}| \times |\mathcal{X}|$ matrix $\mathbf{T}$ of transition
probabilities is \emph{row-stochastic}: $T_{ij} \geq 0$ and
$\sum_j T_{ij} = 1$ for all~$i$.

In HoneyIQ, the attacker's behaviour is modelled with two parallel DTMCs:
%
\begin{itemize}
  \item $\mathbf{T}_A$ ($10 \times 10$): attack-type transitions, governing
        movement between the ten attack categories.
  \item $\mathbf{T}_K$ ($7 \times 7$): kill-chain-stage transitions,
        governing progression through the seven kill chain stages.
\end{itemize}

Intent-specific modifier matrices $\mathbf{M}_A^{(\pi)}$ and
$\mathbf{M}_K^{(\pi)}$ are applied element-wise to the base matrices before
row-renormalisation:
%
\begin{equation}
  \tilde{T}^{(\pi)}_{ij} = T_{ij} \cdot M^{(\pi)}_{ij}, \qquad
  T^{(\pi)}_{ij} = \frac{\tilde{T}^{(\pi)}_{ij}}{\sum_k \tilde{T}^{(\pi)}_{ik}}.
\end{equation}
%
This design separates the \emph{shared attack grammar} (encoded in the base
matrices) from the \emph{intent-specific tendency} (encoded in the modifiers),
allowing four qualitatively different attack campaigns to arise from the
same underlying model.

The SEDM exploits the Markov chain structure directly: the escalation risk
at stage $k$ under intent $\pi$ is:
%
\begin{equation}
  \rho(k, \pi) = \sum_{k' > k} T^{(\pi)}_{kk'},
  \label{eq:esc_risk}
\end{equation}
%
the probability that the attacker advances to a strictly higher kill-chain
stage in the next step.
This quantity is computed analytically from the transition matrix rather than
estimated from trajectory data, making the SEDM's risk assessment exact under
the model.
Because $\rho(k, \pi)$ has no dependence on observed run-time data --- it is
a fixed property of the hand-specified transition matrices --- Chapter~\ref{ch:methodology}
explicitly does not make it a target of the dynamic-adjustment mechanisms
introduced in \S\ref{sec:dynamic_response}; only signals that are genuinely
observational are candidates for runtime adaptation, a distinction made
precise in \S\ref{sec:dynamic_response} and revisited in Section~5.3.

%% ── 2.4 Cyber Kill Chain ─────────────────────────────────────
\section{Lockheed Martin Cyber Kill Chain}
\label{sec:bg_killchain}

The Cyber Kill Chain~\citep{hutchins2011intelligence} models an adversary's
campaign as a sequential progression through seven phases
(Table~\ref{tab:killchain}).
Originally derived from military targeting doctrine, the framework has become
a standard reference model for structuring threat intelligence and mapping
defensive countermeasures.

\begin{table}[ht]
  \centering
  \caption{Cyber Kill Chain stages, associated attack types in HoneyIQ,
           and kill-chain weight values.}
  \label{tab:killchain}
  \begin{tabular}{cllcl}
    \toprule
    Index & Stage & Weight & Primary Attack Types \\
    \midrule
    0 & Reconnaissance    & 0.10 & NORMAL, RECONNAISSANCE \\
    1 & Weaponization     & 0.20 & ANALYSIS \\
    2 & Delivery          & 0.35 & FUZZERS, GENERIC \\
    3 & Exploitation      & 0.55 & EXPLOITS, SHELLCODE \\
    4 & Installation      & 0.70 & BACKDOORS \\
    5 & Command \& Control & 0.85 & WORMS \\
    6 & Actions on Obj.   & 1.00 & DOS \\
    \bottomrule
  \end{tabular}
\end{table}

\paragraph{Reconnaissance (Stage 0).}
The attacker observes the target network: port scanning, service enumeration,
and open-source intelligence (OSINT) gathering.
In HoneyIQ, this maps to RECONNAISSANCE attack type with short-duration,
high-connection-count traffic patterns.
The stage weight of 0.10 reflects that reconnaissance alone does not cause
direct harm.

\paragraph{Weaponization (Stage 1).}
The attacker crafts an exploit or malware tailored to discovered
vulnerabilities.
ANALYSIS traffic (e.g., fuzzing sub-threshold probes) characterises this
phase.

\paragraph{Delivery (Stage 2).}
The weapon is transmitted to the target via phishing, drive-by download,
or network exploitation.
FUZZERS and GENERIC attack types with elevated packet rates represent
delivery-phase traffic in HoneyIQ.

\paragraph{Exploitation (Stage 3).}
The delivered exploit executes code on the target system.
EXPLOITS and SHELLCODE attacks, characterised by high severity and specific
connection patterns, represent this phase.
The stage weight (0.55) reflects the significant danger of active code
execution.

\paragraph{Installation (Stage 4).}
A persistent mechanism (backdoor, rootkit) is installed to maintain access.
BACKDOORS in HoneyIQ exhibit long-duration, low-bandwidth connections
characteristic of covert persistence channels.

\paragraph{Command and Control (Stage 5).}
The attacker establishes a command channel to issue instructions and
exfiltrate data.
WORMS represent C2 activity in HoneyIQ, with high connection counts to
multiple destinations.

\paragraph{Actions on Objectives (Stage 6).}
The attacker executes the final mission objective: data exfiltration,
destruction, or disruption.
DOS attacks with extreme packet rates characterise this stage.
The stage weight of 1.00 reflects maximum threat.

\paragraph{Kill Chain Constraint.}
The Lockheed Martin Kill Chain is a forward-progressing model.
HoneyIQ enforces a floor constraint to prevent unrealistic stage regression:
the sampled next stage is bounded below by the attack type's primary stage
minus one.
This models the attacker's ability to re-use lower-stage techniques while
ensuring they cannot regress arbitrarily.

\paragraph{MITRE ATT\&CK.}
The MITRE ATT\&CK framework~\citep{strom2018mitre} provides a complementary
view at a finer granularity, cataloguing specific techniques within each
kill chain phase.
HoneyIQ's attack-type taxonomy maps to MITRE ATT\&CK at the category level:
RECONNAISSANCE corresponds to the Discovery tactic; EXPLOITATION to the
Execution and Privilege Escalation tactics; WORMS to the Command and
Control tactic.

%% ── 2.5 Honeypot Technology ──────────────────────────────────
\section{Honeypot Technology}
\label{sec:bg_honeypot}

A honeypot is an intentionally vulnerable decoy system designed to attract
attackers and study their behaviour without risk to production
assets~\citep{spitzner2003honeypot, cheswick1992evening}.

\subsection{Taxonomy}

Honeypots are commonly classified along two dimensions:

\paragraph{Interaction level.}
\emph{Low-interaction} honeypots emulate limited services (e.g., an SSH
banner) with minimal actual functionality, reducing risk but limiting
intelligence value.
\emph{High-interaction} honeypots run full operating systems and services,
capturing rich attacker behaviour but requiring careful isolation to prevent
compromise of production infrastructure.

\paragraph{Deployment purpose.}
\emph{Research} honeypots aim to understand attacker techniques and gather
threat intelligence.
\emph{Production} honeypots are deployed within live networks to detect
intrusions and divert attackers.

\subsection{Tarpitting and Deception}

A distinctive honeypot capability is the ability to actively deceive the
attacker rather than passively observe~\citep{provos2004virtual}.
Tarpitting (the TROLL action in HoneyIQ) involves:
%
\begin{itemize}
  \item Responding with synthetic, plausible-but-incorrect data (fake
        credentials, bogus file systems, fabricated network topology).
  \item Introducing artificial delays to waste attacker time and resources.
  \item Feeding false intelligence that misdirects subsequent attack steps.
\end{itemize}
%
Tarpitting is most valuable against persistent attack types (backdoors,
worms, shellcode) where the attacker maintains a long-duration session and
can be engaged for extended observation.

\subsection{Response Actions in HoneyIQ}

HoneyIQ abstracts honeypot management into five actions that span a spectrum
from passive to aggressive:

\begin{enumerate}
  \item \textbf{ALLOW}: pass traffic without logging or intervention.
        Appropriate for confirmed benign traffic to avoid false alarms.
  \item \textbf{LOG}: record the session for later analysis without alerting
        the attacker.
        Appropriate for low-severity, early-stage reconnaissance.
  \item \textbf{TROLL}: engage with fake responses to maximise intelligence
        while wasting attacker resources.
        Appropriate for medium-severity attacks where session duration is
        valuable.
  \item \textbf{BLOCK}: terminate the connection via firewall rules.
        Appropriate for high-severity attacks that should not be permitted
        to continue.
  \item \textbf{ALERT}: issue an immediate high-priority security alert for
        human analyst review.
        Appropriate for critical-severity attacks at late kill-chain stages.
\end{enumerate}

\subsection{Game-Theoretic Perspectives}

The honeypot management problem has been studied from a game-theoretic
perspective~\citep{shi2021learning}, modelling it as a Stackelberg or
Nash equilibrium problem between attacker and defender.
These approaches provide formal optimality guarantees but require explicit
modelling of the attacker's utility function, which is typically unknown in
practice.
RL and decision-matrix approaches avoid this requirement by learning or
designing responses from simulated interactions.

%% ── 2.6 UNSW-NB15 ────────────────────────────────────────────
\section{UNSW-NB15 Dataset}
\label{sec:bg_unsw}

The UNSW-NB15 benchmark~\citep{moustafa2015unsw} was created at the
Australian Centre for Cyber Security to overcome the limitations of the
widely-criticised KDD Cup 99 dataset~\citep{tavallaee2009detailed}.
It contains network flow records captured from a testbed environment running
nine attack categories (Fuzzers, Analysis, Backdoors, DoS, Exploits, Generic,
Reconnaissance, Shellcode, Worms) alongside normal traffic, totalling
approximately 2.5 million records with 49 features.

\subsection{Feature Set}

The 49 UNSW-NB15 features cover five categories:
%
\begin{itemize}
  \item \textbf{Flow-level features}: duration (\texttt{dur}), bytes sent/
        received (\texttt{sbytes}, \texttt{dbytes}), packets sent/received
        (\texttt{spkts}, \texttt{dpkts}), load in bits/s
        (\texttt{sload}, \texttt{dload}).
  \item \textbf{IP-level features}: time-to-live values
        (\texttt{sttl}, \texttt{dttl}), packet loss counts
        (\texttt{sloss}, \texttt{dloss}).
  \item \textbf{TCP-level features}: window sizes (\texttt{swin}, \texttt{dwin}).
  \item \textbf{Connection-level features}: \texttt{ct\_srv\_src}
        (connections to the same service from the same source),
        \texttt{ct\_dst\_ltm} (connections to the same destination in the
        last time interval).
  \item \textbf{Content features}: protocol, service, state.
\end{itemize}

HoneyIQ uses 15 of these features, selected for their discriminatory power
across attack categories.
Rather than fitting distributions to the actual UNSW-NB15 records, the
simulation samples from \emph{parametric approximations} that reproduce the
qualitative signatures of each attack type (e.g., high \texttt{sload} for
DoS, long \texttt{dur} for Backdoors).
\S\ref{sec:realism} extends this generator with a session-level intensity
multiplier and, for benign traffic, a choice among three behavioural
personas, so that repeated samples from the same simulated session or the
same benign source are no longer independent draws.

\subsection{Attack Signatures in HoneyIQ}

Table~\ref{tab:feature_signatures} summarises the distinctive per-attack
feature patterns implemented in HoneyIQ.

\begin{table}[ht]
  \centering
  \caption{Characteristic UNSW-NB15 feature patterns per attack type in HoneyIQ.}
  \label{tab:feature_signatures}
  \begin{tabular}{lll}
    \toprule
    Attack Type & Distinctive Feature Pattern \\
    \midrule
    NORMAL         & Short \texttt{dur}, moderate balanced \texttt{sbytes/dbytes} \\
    RECONNAISSANCE & High \texttt{ct\_dst\_ltm} ($\lambda=30$), many short connections \\
    ANALYSIS       & Moderate \texttt{sload}, medium \texttt{ct\_srv\_src} \\
    FUZZERS        & High \texttt{spkts} (Poisson $\lambda=100$), varied payloads \\
    GENERIC        & Balanced traffic, similar to normal but higher packet counts \\
    EXPLOITS       & High \texttt{sbytes}, spike in \texttt{sload} (10k--100k bps) \\
    SHELLCODE      & Low \texttt{dur}, very high \texttt{sload}, small payloads \\
    BACKDOORS      & Long \texttt{dur} (10--3600\,s), low stealthy \texttt{sload} \\
    DOS            & Extreme \texttt{sload} (50k--500k bps), near-zero \texttt{dur} \\
    WORMS          & High \texttt{ct\_dst\_ltm} ($\lambda=40$), spreading pattern \\
    \bottomrule
  \end{tabular}
\end{table}

%% ── 2.7 Random Forests ───────────────────────────────────────
\section{Random Forests}
\label{sec:bg_rf}

A Random Forest~\citep{breiman2001random} is an ensemble of decision trees,
each trained on a bootstrap sample of the training data with a randomly
selected subset of features considered at each split.
Aggregating predictions across many uncorrelated trees reduces variance
without increasing bias, yielding classifiers that are robust to overfitting
and effective on high-dimensional data with mixed feature types.

In HoneyIQ, the Random Forest classifier serves as a supporting component
for attack-type identification:
%
\begin{itemize}
  \item \textbf{Training data}: 6,000 synthetic samples (600 per class)
        generated from the parametric feature distributions.
  \item \textbf{Class balance}: \texttt{class\_weight=`balanced'} reweights
        the contribution of each sample inversely proportional to its class
        frequency, preventing the majority class from dominating gradients.
  \item \textbf{Feature normalisation}: \texttt{StandardScaler} zero-centres
        and unit-scales each feature before training and inference.
  \item \textbf{Configuration}: 150 trees, maximum depth 20.
\end{itemize}

Because training data is generated directly from the feature distributions,
the classifier has an inherent advantage over real-world settings: there is
no distribution shift between training and evaluation data.
High accuracy ($>$\,95\% on held-out synthetic data) is therefore expected,
and the classifier primarily serves as a preprocessing step for logging and
introspection rather than as the primary decision-making mechanism.

\begin{quote}
\noindent\textit{Note (corrected results, Chapter~\ref{ch:results}):} the
held-out accuracy of the classifier trained on the original
(pre-realism-extension) generator is 99.85\% (10-class), well above the
``$>$\,95\%'' figure anticipated in this background chapter, and the SEDM's
evaluation now runs in a classifier-driven decision mode rather than
treating the classifier purely as a logging/introspection component --- see
\S\ref{sec:classifier} and Chapter~\ref{ch:results}. Retraining the
classifier on the session-coherent, persona-diverse generator introduced in
\S\ref{sec:realism} changes this figure; the updated value and its effect on
the classifier-driven false-positive rate are reported in Section~4.9.
\end{quote}

%% ── 2.8 Interpretable Machine Learning ───────────────────────
\section{Interpretable Machine Learning}
\label{sec:bg_interpretable}

The past decade has seen rapid progress in predictive performance of neural
networks at the cost of reduced interpretability.
In high-stakes domains such as healthcare, criminal justice, and
cybersecurity, this trade-off is often unacceptable~\citep{rudin2019stop}.

\subsection{Motivation for Interpretability in Security}

Cybersecurity deployments have specific interpretability requirements:
%
\begin{itemize}
  \item \textbf{Auditability}: security operations centres (SOCs) must be
        able to explain to auditors and regulators why a particular IP was
        blocked or why an alert was raised.
  \item \textbf{Trust and override}: analysts must be able to identify when
        a policy is making systematic errors and override its decisions.
  \item \textbf{Debugging}: when a policy allows a known attack (false
        negative) or blocks benign traffic (false positive), the responsible
        decision rule should be identifiable and correctable.
  \item \textbf{Adversarial robustness}: opaque policies are potentially
        vulnerable to adversarial inputs crafted to manipulate neural network
        activations; interpretable policies are more resistant because their
        logic is explicit.
\end{itemize}

\subsection{Decision Tables as Interpretable Policies}

A decision table maps a discretisation of the input space to actions.
It is interpretable by construction: any decision can be traced to a
specific cell in the table and the associated input conditions.
Decision tables have been used in medical diagnosis, industrial control
systems, and credit scoring.
In the security domain, they underpin classic IDS rule engines such as Snort.

The SEDM in HoneyIQ is a $7 \times 3$ decision table (stage $\times$ band)
with, in this version, four supplementary override rules
(\S\ref{sec:sedm}).
This structure is small enough to be printed on a single page, audited by
a practitioner in minutes, and modified without retraining.

\subsection{Trade-offs}

Decision tables achieve interpretability by discretising a continuous space.
This discretisation introduces step discontinuities at band boundaries:
two states with escalation risks of 0.349 and 0.351 receive different base
actions (Low vs.\ Medium band), even though the underlying risk is nearly
identical.
The override rules partially mitigate this by providing additional conditions
that can upgrade the action regardless of the band boundary.
Chapter~\ref{ch:discussion} examines this trade-off empirically, and returns
to it specifically when discussing why the dynamic-response mechanisms in
\S\ref{sec:dynamic_response} were deliberately kept in this same
interpretable-by-construction family rather than delegated to a learned
adjustment layer.

%% ── 2.9 Related Work ─────────────────────────────────────────
\section{Related Work}
\label{sec:bg_related}

\subsection{RL-Based Intrusion Response}

Several studies have applied reinforcement learning to network intrusion
response.
\citet{malialis2015distributed} applied Q-learning agents to distributed
denial-of-service (DDoS) mitigation, demonstrating that independent
per-router agents could collectively throttle attack traffic without
centralised coordination.
\citet{elderman2017adversarial} studied adversarial RL in a network
intrusion simulation, showing that a defender agent trained against a fixed
attacker degrades when the attacker adapts.
\citet{hammar2021learning} framed intrusion prevention as an optimal stopping
problem, using RL to determine when to block a suspicious connection.
\citet{li2019mad} combined multi-agent RL with inverse reinforcement learning
to infer attacker utility functions and design countervailing responses.

HoneyIQ extends this body of work by (1) modelling attacker intent as a
latent variable that governs the Markov transition structure, enabling
evaluation of policy robustness across qualitatively distinct campaigns, and
(2) introducing the SEDM as an interpretable alternative to learned policies
that achieves comparable detection performance, and (3), in this version, by
directly confronting the question these RL-based approaches implicitly
answer in the affirmative --- whether a learned policy is the appropriate
mechanism for \emph{adaptive} behaviour --- with project-specific evidence
rather than assuming it (Section~5.3).

\subsection{Attacker Simulation}

Realistic attacker simulation is a prerequisite for evaluating any defensive
policy~\citep{applebaum2017cyber}.
Agent-based simulators such as CyberBattle Gym and CAGE Challenge environments
model network topologies and service vulnerabilities explicitly.
HoneyIQ abstracts topology to focus on the temporal dynamics of attack
progression, using the Markov chain to produce plausible kill-chain sequences
without requiring a specific network graph.
The UNSW-NB15-inspired feature distributions provide a bridge to real-world
traffic statistics while remaining computationally tractable;
\S\ref{sec:realism} extends this bridge with session-level coherence and
multi-persona benign traffic, addressing the specific gap that a
single-shot, i.i.d.-per-step feature simulator does not model a session as
originating from one consistent source.

\subsection{Honeypot Optimisation}

\citet{shi2021learning} modelled honeypot placement as a bilevel optimisation
problem and used RL to find strategies that maximise the attacker's deception
cost.
\citet{franco2021survey} provided a comprehensive survey of honeypot
architectures and their deployment in production environments.
The HoneyIQ action abstraction (ALLOW / LOG / TROLL / BLOCK / ALERT) captures
the essential operational decisions without committing to a specific honeypot
technology, making the framework applicable to a range of deployment contexts.

\subsection{Partially Observable Defence}

\citet{miehling2018optimal} studied optimal defence under partial
observability using a POMDP framework with Bayesian attack graphs.
HoneyIQ's primary and extended evaluations (Chapter~\ref{ch:results}) report
both an oracle mode, treating the attacker's state (attack type, kill chain
stage, intent) as fully observable via the ground-truth state vector, and a
classifier-driven mode that decides from the Random Forest classifier's
predicted attack type instead, introducing a bounded form of partial
observability.
Re-evaluating this classifier-driven mode under realistic (non-synthetic)
classification noise remains a priority direction for future work
(Section~\ref{sec:future}).

\begin{quote}
\noindent\textit{Editorial note on this chapter:} two passages in the
original chapter (\S2.6, \S2.9's closing paragraph) described the
classifier's role and accuracy using language written before the
classifier-driven evaluation mode existed and before the corrected
evaluation pipeline was re-run (see Chapter~\ref{ch:results} and
\texttt{docs/BUGS\_AND\_FIXES.md}, Bug 8/9). The notes above adjust the
reading of those passages for consistency with the corrected results
without rewriting the original prose; see
\texttt{docs/parameter\_selection.md} and
\texttt{docs01/parameter\_selection.md} for full technical detail.
\end{quote}

%% ── 2.10 Reputation-Based Access Control and Bounded Control ─
\section{Reputation-Based Access Control and Bounded Control}
\label{sec:bg_repctrl}

The two dynamic-response mechanisms introduced in \S\ref{sec:dynamic_response}
each draw on a distinct, well-established design pattern from outside the
reinforcement-learning literature reviewed above; this section introduces
both briefly.

\subsection{Reputation and Escalating Response}

Maintaining a persistent, decaying trust or offense score per network
identity --- and escalating the response once that score crosses a
threshold, independent of how any single subsequent event looks in
isolation --- is a long-standing pattern in operational network defence.
Widely-deployed tools such as fail2ban implement exactly this idea at the
level of a single host: repeated authentication failures from one source
accumulate into a ban, and the ban persists across the individual connection
attempts that triggered it rather than resetting with each new attempt.
The pattern generalises beyond any one tool: what matters structurally is
(1) a score that survives longer than the individual session or connection
that updates it, (2) a decay function so that old behaviour eventually stops
mattering, and (3) a threshold-driven escalation that can override what a
single, isolated observation would otherwise suggest.
HoneyIQ's \texttt{ReputationTracker} (\S\ref{sec:dynamic_response}) follows
this structure directly, using a continuous half-life decay in place of a
fixed ban duration, and integrates it into the SEDM as an explicit, numbered
override rule (R4) rather than a separate blocklist mechanism, so that a
reputation-driven decision remains traceable through the same audit path as
every other SEDM decision.

\subsection{Bounded (Deadband) Control}

A deadband controller is a standard control-theoretic pattern for keeping a
measured quantity near a target without responding to noise: the controller
takes no action while the measured value is within a tolerance band around
the target, and nudges a control variable in the corrective direction only
once the measurement drifts outside that band --- typically with the control
variable itself clamped to a bounded range to limit how far the system can
be pushed from its default operating point.
This is a substantially weaker (and more auditable) mechanism than a learned
controller: it requires no training signal, its behaviour is fully
determined by a handful of named constants (target, tolerance, step size,
bound), and its entire state at any moment is a single number that can be
read and explained directly.
HoneyIQ's \texttt{AdaptiveThresholds} (\S\ref{sec:dynamic_response}) is
precisely this pattern, applied to the SEDM's R3 trigger-rate rather than to
a physical process variable --- a deliberate choice, discussed in
\S\ref{sec:dynamic_response} and revisited in Section~5.3, to use the
\emph{least} mechanistically complex tool that could plausibly address the
alert-fatigue problem it targets, rather than defaulting to a learned
adjustment mechanism.
